\newpage
\section{Vorgehensweise und Methodik}

Die vorliegende Arbeit ist theoretisch-anwendungsorientiert angelegt. Ziel ist die Entwicklung
eines Artefakts im Sinne der Wirtschaftsinformatikforschung, konkret eines Vorgehensmodells
für den Einsatz von UI-basierter Robotic Process Automation in API-losen Legacy-Umgebungen.
Als methodischer Rahmen soll daher Design Science Research (DSR) dienen, da dieser Ansatz
explizit auf die gestaltungsorientierte Erzeugung von Wissen durch die Entwicklung nützlicher
Artefakte ausgerichtet ist \autocite[Vgl.][S.~75]{Hevner2004}. Die Entscheidung für DSR ergibt sich aus dem
Charakter der Forschungsfrage: Es soll kein bestehendes Phänomen beschrieben, sondern ein
strukturiertes Hilfsmittel für die Praxis konstruiert werden. Diese Zielsetzung lässt sich im
Methodenspektrum der Wirtschaftsinformatik dem gestaltungsorientierten Paradigma zuordnen,
das neben dem verhaltenswissenschaftlichen Paradigma als eigenständiger Forschungsansatz
anerkannt ist \autocite[Vgl.][S.~280]{WildeHess2007}.

Innerhalb des DSR-Rahmens soll die Arbeit dem sechsphasigen Prozessmodell nach
\textcite{Peffers2007} folgen, das von der Problemidentifikation über die Zieldefinition und
Artefaktentwicklung bis hin zur Evaluation und Kommunikation der Ergebnisse reicht. Dieses
Modell eignet sich insbesondere für IS-Artefakte vom Typ Methode oder Modell \autocite[Vgl.][S.~253]{MarchSmith1995}
und ermöglicht eine nachvollziehbare Strukturierung des gesamten Forschungsprozesses. Die
drei Forschungszyklen nach \textcite{Hevner2007} -- Relevanz-, Rigor- und Designzyklus --
sollen dabei als übergeordnetes Qualitätskriterium herangezogen werden, um sicherzustellen,
dass das entwickelte Modell sowohl wissenschaftlich fundiert als auch praxisrelevant ist.

\newpage
\section{Vorgehensweise und Methodik}

% TODO: Gewaehlte Methode benennen und begruenden.

% TODO: Datenerhebung beschreiben.

% TODO: Auswertungsverfahren / Theorien beschreiben.

% --- GRAFIK: DSRM-Prozessmodell (Peffers et al. 2007) ---
\begin{figure}[H]
  \caption{Design Science Research Methodology: Prozessmodell 
           \autocite[Vgl.][54]{Peffers2007}}
  \label{fig:dsrm_prozess}
  \centering
  \includegraphics[angle=270, origin=c, width=0.90\textwidth]{peffers_dsrm}
\end{figure}
% ---------------------------------------------------------

Die Grundlage für das zu entwickelnde Vorgehensmodell soll durch eine strukturierte
Literaturrecherche gelegt werden. Diese soll dem Vorgehen nach \textcite{VomBrocke2009}
folgen, das eine transparente Dokumentation der Suchstrategie, der verwendeten Datenbanken
und der Einschluss- sowie Ausschlusskriterien verlangt. Geplant ist die Suche in den
einschlägigen Datenbanken der Wirtschaftsinformatik und der Informatik, darunter SpringerLink,
ACM Digital Library und IEEE Xplore. Die Recherche soll die Themenbereiche RPA-Grundlagen,
UI-Automatisierung, Legacy-Systemintegration sowie bestehende Vorgehensmodelle für
Automatisierungsvorhaben abdecken. Auf Basis dieser Literatur sollen vorhandene Modelle und
Frameworks analysiert und hinsichtlich ihrer Übertragbarkeit auf den API-losen Kontext
bewertet werden.

Die eigentliche Artefaktentwicklung soll iterativ erfolgen: Ein erster Modellentwurf wird aus
der Literaturanalyse abgeleitet und anschließend einer Evaluation unterzogen. Als
Evaluationsstrategie ist -- im Sinne des FEDS-Frameworks \autocite[Vgl.][S.~77]{VenableEtAl2016} --
eine formative Expertenvalidierung vorgesehen. Dazu sollen leitfadengestützte Gespräche mit
Praktikern geführt werden, die über nachweisliche Erfahrung in der Einführung oder dem Betrieb
von RPA-Lösungen in Legacy-Umgebungen verfügen.
Der Stichprobenumfang und die Auswahlkriterien für die Expertinnen und Experten werden im
Rahmen der Thesis methodisch begründet; zum jetzigen Zeitpunkt ist ein Umfang von fünf bis
acht Gesprächen als angemessen einzuschätzen, um inhaltliche Sättigung zu erreichen, ohne
den Rahmen einer Bachelorarbeit zu überschreiten. Die Erkenntnisse aus der Expertenvalidierung
sollen abschließend in eine überarbeitete Modellversion einfließen, sodass ein iterativ
verfeinertes und empirisch rückgekoppeltes Vorgehensmodell als Ergebnis der Arbeit vorliegt.

Hinsichtlich der wissenschaftlichen Positionierung soll die Arbeit dem Anspruch folgen,
ein neuartiges Artefakt für einen bisher unzureichend adressierten Anwendungskontext zu
entwickeln \autocite[Vgl.][S.~337]{GregorHevner2013}. Im Rahmen dieser Arbeit wird daher die These
vertreten, dass bestehende RPA-Einführungsmodelle den spezifischen Anforderungen
API-loser Legacy-Umgebungen nicht hinreichend gerecht werden und ein kontextspezifisches
Vorgehensmodell einen eigenständigen wissenschaftlichen und praktischen Beitrag darstellt.





% Leitfragen (FOM-Gliederungshilfe):
% - Theoretisch oder theoretisch-anwendungsorientiert? Eigene empirische
%   Untersuchung wird bevorzugt. Moegliche Ansaetze:
%     * Literatur-Review (vgl. vom Brocke u.a.)
%     * Design Science (vgl. Peffers et al. / Hevner et al.)
%     * qualitative Studie (vgl. Mayring, Kuckartz, ...)
%     * quantitative Untersuchung
%   (Mixed-Methods ist fuer eine Thesis i.d.R. zu gross.)
% - Eigene Datenerhebung oder Stuetzung auf vorhandene Daten / Literatur?
% - Erhebungsmethode (Fragebogen, Interviews, Expertengespraeche, ...)? Bei wem?
%   Stichprobenumfang?
% - Wie werten Sie die Daten aus? Welche Theorie(n) setzen Sie ein?

% TODO: Gewaehlte Methode benennen und begruenden.

% TODO: Datenerhebung beschreiben.

% TODO: Auswertungsverfahren / Theorien beschreiben.
